% !TEX program = lualatex
% =====================================================================
% execucao-cmo.tex -- nota tecnica
%
% Previsao do custo marginal de operacao semi-horario contra a ingenua
% sazonal, e o que a distribuicao do alvo impoe a comparacao.
%
% Todo numero vem de data/interim/E-preco-cmo/: janelas.json, controles.json,
% resultado.json, decomposicao.json e ablacao.json, lidos do arquivo e nao de
% saida de tela. Onde um numero aparece aqui e no README do experimento, os
% dois saem do mesmo JSON -- mudar um exige mudar o outro, porque a ligacao e
% de procedencia e nao automatica.
%
% Nao entra no `reporte.tex`. Compila sozinho:
%   latexmk -pdf execucao-cmo.tex
% =====================================================================

% A folha de estilo pressupoe `book`: usa `\chaptertitlename` e os
% comandos de capitulo do `tocloft`. Em `article` ela nao carrega.
\documentclass[10pt,a4paper,oneside,twocolumn]{book}

\usepackage{iftex}
\ifPDFTeX\usepackage[utf8]{inputenc}\fi
\usepackage[brazilian]{babel}
% Mancha e medianiz do `reporte.tex`: 17,2 cm em duas colunas de 8,25 cm.
\usepackage[a4paper,top=2.0cm,bottom=2.2cm,left=1.9cm,right=1.9cm,
            columnsep=0.7cm]{geometry}
\usepackage{estilo-reporte}

\raggedbottom
\pagestyle{plain}
\setlength{\parskip}{0.35em}
\setlength{\parindent}{0pt}

\marca{../docs/img/terra-mark.png}
\autoria{João Leonardi da Silva Melo}{joao\_leonardi.melo@somosicev.com}

\begin{document}
\abertura{Previsão do custo marginal de operação}
         {Nota técnica · \today{} ·
          \texttt{experiments/E-preco-cmo/}}

\section*{Objetivo}

Medir a habilidade de modelos de aprendizado de máquina sobre a linha de base
ingênua sazonal na previsão do CMO semi-horário do SIN, a um e a dois dias, e
determinar o que na distribuição do alvo governa o resultado.

Tarefa~11 de \texttt{tarefas-dados.tex}. Até 6 de setembro de 2026 a tarefa
tinha dado em disco e nenhum experimento: \texttt{data/raw/cmo/} constava do
inventário como bruto órfão, baixado e não lido por código.

\section*{Dados}

CMO semi-horário do ONS, três arquivos anuais. Covariáveis: curva de carga
horária, intercâmbio nacional horário, EAR e ENA diárias por subsistema.

\begin{ficha}
  \item[Fonte]     \num{184704} registros, 4 subsistemas $\times$ \num{46176}
                   meias-horas, 01/01/2024 a 27/08/2026, zero nulos.
  \item[Painel]    \num{180288} linhas de \num{186240} na grade cheia.
  \item[Descartes] \num{1536} sem alvo (8 dias de CMO ausentes na fonte) e
                   \num{4416} sem D$-$1 ou D$-$7.
  \item[Alvo]      CMO em R\$/MWh da meia-hora, mínimo $-39{,}24$,
                   máximo \num{4870.95}.
\end{ficha}

\campo{Massa no piso} 20,5\% das meias-horas do painel ficam em CMO
$\leq 1$~R\$/MWh. A distribuição é inflada no piso e de cauda pesada à direita,
com p99 entre 614 e \num{1082} conforme o subsistema.

\campo{Cada buraco na fonte custa três dias de painel} Além do próprio dia, cai
o dia seguinte, que perde o D$-$1 da ingênua, e o dia sete depois, que perde o
D$-$7. Oito buracos, \textbf{24 dias fora do painel}. Os dias ausentes não são
interpolados: interpolar criaria alvo que o operador não publicou.

\section*{Método}

\campo{Estimando} CMO das 48 meias-horas do dia D nos quatro subsistemas, com
previsão emitida às 12:00 do dia D$-$1.

\campo{Adversário} Ingênua sazonal: terça a sexta repetem D$-$1; sábado,
domingo e segunda repetem D$-$7.

\campo{Perda} MAE, declarada antes da execução. RMSE como secundário.

\campo{Origem móvel} Recalibração mensal com janela de treino expansiva, 15
reajustes. Teste de 01/06/2025 a 27/08/2026: \num{85248} linhas, 444 dias,
41 features.

\campo{Intervalo} Reamostragem de dia, \num{2000} repetições, os quatro
subsistemas no mesmo bloco.

\vspace{0.2em}
\tabfont
\begin{tabular}{@{}>{\rag\arraybackslash}p{3.0cm}>{\rag\arraybackslash}p{4.6cm}@{}}
\toprule
\cab{Fonte} & \cab{Portão} \\
\midrule
CMO                 & até o fim de D$-$1 \\
Carga, intercâmbio  & até o fim de D$-$2 \\
EAR, ENA            & até D$-$3 \\
Calendário de D     & determinístico \\
\bottomrule
\end{tabular}
\notas{Portão medido, uma observação por fonte, pela diferença entre o último
registro do arquivo e a data de modificação dele. O CMO é publicado
\textbf{antes} do dia: o arquivo baixado às 10:08 de 27/08/2026 já trazia as 48
meias-horas daquele mesmo dia. Carga e intercâmbio atrasam 1,2~d; EAR e ENA,
2,2 e 2,3~d. O portão vale igual para o modelo e para as linhas de base.}

\campo{Modelos} Três variantes de LightGBM sobre o mesmo insumo: \emph{direto}
prevê o nível do preço; \emph{residual} prevê a correção da ingênua;
\emph{duas partes} classifica se a meia-hora cai no piso e regride o resto, com
a regra de decisão que minimiza MAE sob esse modelo. Uma configuração de
hiperparâmetro, uma semente.

\section*{Resultados}

\vspace{0.2em}
\tabfont
% A coluna de IC leva $[-176{,}4;-136{,}1]$, que estoura 1,75 cm com o
% `\tabcolsep` padrao. Reduzir a separacao paga a largura sem baixar o corpo.
\begingroup\setlength{\tabcolsep}{3pt}
\begin{tabular}{@{}>{\rag\arraybackslash}p{2.05cm}rrr>{\rag\arraybackslash}p{2.05cm}@{}}
\toprule
\cab{Preditor} & \cab{MAE} & \cab{RMSE} & \cab{Hab.} & \cab{IC 95\%} \\
\midrule
Ingênua D$-$1    &  52,70 & 129,33 & $-7{,}36$   & $[-14{,}1;-1{,}2]$ \\
Ingênua D$-$7    &  69,82 & 140,19 & $-42{,}24$  & $[-53{,}7;-32{,}0]$ \\
\textbf{Ingênua sazonal} & \textbf{49,09} & \textbf{122,02} & --- & --- \\
Climatologia     & 125,18 & 162,89 & $-155{,}02$ & $[-176{,}4;-136{,}1]$ \\
LGBM direto      &  51,32 & 103,38 & $-4{,}54$   & $[-12{,}4;+2{,}7]$ \\
LGBM residual    &  50,04 & 115,42 & $-1{,}94$   & $[-6{,}5;+2{,}3]$ \\
\textbf{LGBM duas partes} & \textbf{48,71} & \textbf{102,31} & $\mathbf{+0{,}76}$ & $[-6{,}7;+7{,}7]$ \\
Oráculo de nível &  40,42 & 112,69 & $+17{,}66$  & $[+15{,}3;+20{,}2]$ \\
\bottomrule
\end{tabular}
\endgroup
\notas{MAE e RMSE em R\$/MWh. Habilidade em \% sobre a ingênua sazonal, na
métrica MAE. Habilidade em RMSE: duas partes $+16{,}15$, direto $+15{,}27$,
residual $+5{,}41$. Climatologia é a média por subsistema e meia-hora,
calculada no treino de cada dobra. O oráculo recebe a mediana verdadeira do
resíduo do dia: não é previsão, é controle positivo.}

\vspace{0.4em}
Os três modelos têm intervalo cruzando zero em MAE. Em RMSE os três batem a
ingênua, o de duas partes por 16,2\%.

\campo{Horizonte} O alvo deslocado para D$+$1, com o portão inteiro deslocado
junto.

\vspace{0.2em}
\tabfont
\begin{tabular}{@{}>{\rag\arraybackslash}p{1.3cm}rrr>{\rag\arraybackslash}p{1.85cm}@{}}
\toprule
\cab{Horiz.} & \cab{Modelo} & \cab{Ingênua} & \cab{Hab.} & \cab{IC 95\%} \\
\midrule
D     & 48,71 & 49,09 & $+0{,}76$  & $[-6{,}7;+7{,}7]$ \\
D$+$1 & 62,68 & 74,89 & $\mathbf{+16{,}31}$ & $[+10{,}6;+21{,}9]$ \\
\bottomrule
\end{tabular}
\notas{MAE em R\$/MWh, modelo de duas partes. Degradação de um dia para o
outro: ingênua $+25{,}80$, modelo $+13{,}96$. O erro cresce nos dois termos, o
que o controle de horizonte exige.}

\campo{Comportamento no piso} 14,0\% das linhas de teste têm o alvo em
CMO $\leq 1$~R\$/MWh.

\vspace{0.2em}
\tabfont
\begin{tabular}{@{}>{\rag\arraybackslash}p{2.5cm}rrr@{}}
\toprule
\cab{Preditor} & \cab{No piso} & \cab{Md.\ ali} & \cab{Md.\ fora} \\
\midrule
Ingênua sazonal  & 71,6\% &  0,00 &  15,10 \\
LGBM direto      & 19,2\% & 16,65 &  22,86 \\
LGBM residual    & 18,6\% & 13,06 &  19,48 \\
LGBM duas partes & 66,0\% &  0,01 &  22,03 \\
Climatologia     &  0,0\% & 97,01 & 127,38 \\
\bottomrule
\end{tabular}
\notas{\emph{No piso}: fração das linhas de alvo no piso em que o preditor
também prevê $\leq 1$~R\$/MWh. \emph{Md.}: erro absoluto mediano, R\$/MWh,
dentro e fora dessas linhas.}

\vspace{0.4em}
A ingênua é um caminho realizado e acerta o piso exatamente. Um regressor único
devolve a mediana condicional, que quase nunca cai em zero. Modelar o piso
explicitamente move a habilidade de $-4{,}54$\% para $+0{,}76$\%.

% Rotulo e tabela nao podem cair em colunas diferentes: o `\campo` sozinho no
% pe da coluna vira titulo orfao. A minipage os mantem juntos.
\begin{minipage}{\linewidth}
\campo{Quantis do erro absoluto}

\vspace{0.2em}
\tabfont
\begin{tabular}{@{}>{\rag\arraybackslash}p{2.15cm}rrrrr@{}}
\toprule
\cab{Preditor} & \cab{q25} & \cab{q50} & \cab{q75} & \cab{q90} & \cab{q95} \\
\midrule
Ingênua sazonal  & \textbf{2,25} & \textbf{12,10} & \textbf{53,50} & 157,6 & 222,5 \\
LGBM direto      & 7,73 & 22,06 & 69,94 & \textbf{136,5} & \textbf{179,9} \\
LGBM duas partes & 5,13 & 19,78 & 63,68 & 131,8 & 184,9 \\
\bottomrule
\end{tabular}
\notas{R\$/MWh. q99: ingênua 346,7; direto 311,3; duas partes 315,7. Fração de
linhas em que o preditor erra menos que a ingênua: direto 39,1\%, residual
41,3\%, duas partes 40,7\%, oráculo 61,1\%.}
\end{minipage}

\vspace{0.4em}
A ingênua erra menos nos quantis baixos e mais nos altos. MAE pesa toda linha
igual; RMSE pesa a cauda.

\campo{Nível do dia e forma dentro do dia}

\vspace{0.2em}
\tabfont
\begin{tabular}{@{}>{\rag\arraybackslash}p{2.5cm}rrr@{}}
\toprule
\cab{Preditor} & \cab{Completo} & \cab{Nível} & \cab{Forma} \\
\midrule
Ingênua sazonal  & 49,09 & 40,21 & 48,64 \\
LGBM direto      & 51,32 & 37,80 & 42,45 \\
LGBM duas partes & \textbf{48,71} & \textbf{35,52} & 42,65 \\
Oráculo de nível & 40,42 & 23,57 & 48,64 \\
\bottomrule
\end{tabular}
\notas{MAE em R\$/MWh. \emph{Nível}: média do dia por subsistema.
\emph{Forma}: desvio de cada meia-hora em torno dessa média. O nível responde
por 39,3\% da variância do CMO no conjunto dos quatro subsistemas e de 32,2\% a
52,4\% dentro de cada um, na janela de teste.}

\vspace{0.4em}
Os modelos batem a ingênua nas duas partes medidas em separado e empatam no
total. MAE não é aditiva sobre a separação: o erro de uma linha é a soma dos
dois, e $|a+b|$ não é $|a|+|b|$.

\campo{Dispersão} Habilidade do modelo de duas partes, por subsistema:
S $+3{,}61$\%, SE $+1{,}49$\%, NE $+0{,}89$\%, N $-2{,}68$\%. Por mês de teste,
de $-68{,}2$\% (dez/2025) a $+38{,}7$\% (jan/2026), acima de zero em 10 dos 15
meses. Dois meses carregam o agregado.

\section*{Controles}

\vspace{0.2em}
\tabfont
\begin{tabular}{@{}>{\rag\arraybackslash}p{0.5cm}>{\rag\arraybackslash}p{3.1cm}>{\rag\arraybackslash}p{3.6cm}@{}}
\toprule
& \cab{Controle} & \cab{Resultado} \\
\midrule
C1 & Grade, chave única e dias ausentes preservados &
     0 duplicatas, 0 dias incompletos, 24 ausentes \\
C2 & Faixa física do alvo &
     0 fora de faixa, $-39{,}24$ a \num{4870.95} \\
C3 & Portão de informação &
     40 dias, \num{307200} células, 0 divergências \\
C4 & Cobertura de feature $\geq$ 97\% &
     pior 99,15\% (\texttt{cmo\_l2}) \\
C5 & Ordem das linhas de base &
     43,93 $<$ 136,40 e 42,47 $\leq$ 61,75 \\
C6 & Subsistemas são séries distintas &
     correlação máxima 0,933 \\
\bottomrule
\end{tabular}
\notas{Os seis são bloqueantes e os seis passaram. O C3 apaga das fontes tudo o
que está depois do instante de emissão e reconstrói as features com a mesma
função do passo~1, exigindo valor idêntico célula a célula. Nulo por permutação
de feature no treino, com alvo intacto: habilidade $-206{,}91$\%.}

\campo{O C1 reprovou uma vez} O limiar estava em 8, o número de buracos na
fonte; o painel perde 24 dias. A aritmética não estava medida.

\campo{O controle positivo custou uma correção} O oráculo deslocado pela média
do resíduo do dia entrega $+0{,}9$\% com intervalo cruzando zero. A perda é MAE
e a constante que a minimiza é a mediana. Trocada, $+17{,}66$\% com intervalo
inteiro acima de zero, sobre o mesmo dado e o mesmo dia.

\section*{Ablação}

\vspace{0.2em}
% Tabela e nota juntas: a nota define `Canon.` e `1a cor.`, e lida uma coluna
% depois da tabela ela nao serve para nada.
\begin{minipage}{\linewidth}
\tabfont
\begin{tabular}{@{}>{\rag\arraybackslash}p{2.55cm}rrrr@{}}
\toprule
\cab{Conjunto} & \cab{$n$} & \cab{Canôn.} & \cab{1ª cor.} & \cab{Dif.} \\
\midrule
A0 calendário + CMO próprio      & 20 & $-4{,}87$ & $-3{,}66$ & 1,21 \\
A1 + CMO dos outros subsistemas  & 28 & $-1{,}01$ & $-1{,}84$ & 0,83 \\
A2 + carga e intercâmbio         & 36 & $-9{,}50$ & $-2{,}69$ & \textbf{6,81} \\
A3 + hidrologia (completo)       & 41 & $\mathbf{+0{,}76}$ & $+4{,}40$ & 3,64 \\
\bottomrule
\end{tabular}
\notas{Habilidade em \%, modelo de duas partes, mesma origem móvel e mesma
recalibração do resultado principal. \emph{Canôn.}: colunas na ordem do painel.
\emph{1ª cor.}: colunas na ordem em que a ablação as concatenava. As duas
corridas usam o mesmo conjunto de features e a mesma semente. \emph{Dif.} em
pontos percentuais.}
\end{minipage}

\vspace{0.4em}
Com \texttt{colsample\_bytree} menor que 1 o LightGBM sorteia coluna por
posição: permutar as mesmas colunas equivale a trocar de semente. Na ordem
canônica o A3 reproduz o resultado principal exatamente, $+0{,}76$\% contra
$+0{,}76$\%, o que fecha a explicação.

\begin{objecao}
As diferenças entre degraus vizinhos --- 3,86, $-8{,}49$ e 10,26 pontos --- são
da mesma ordem que a variação induzida só por permutar coluna, de 0,83 a 6,81.
Com uma semente, a ablação mede variância de ajuste junto com efeito de família
e não separa as duas. \textbf{Nenhuma afirmação sobre qual fonte carrega a
habilidade se sustenta com este número de sementes}, e esta nota não faz
nenhuma. A faixa inteira cabe dentro do intervalo de 14 pontos do resultado
principal.
\end{objecao}

\section*{Limitações}

\campo{O portão é um ponto por fonte} Os arquivos do ONS não trazem carimbo de
publicação. A data de modificação do arquivo em disco é a única evidência
disponível, e o portão inteiro repousa sobre cinco observações.

\campo{O CMO é saída de modelo} É o valor que o modelo de despacho do operador
produz, não medição ex-post. Não há segunda fonte para conferir contra, e
nenhum controle testa se o número publicado é o custo marginal realizado.

\campo{Sem feriado} Não há tabela de feriado no repositório. O efeito cai
dentro do erro dos dois termos, e a ingênua sazonal é a base que mais sofre com
isso.

\campo{Sem preço de liquidação} O CMO não é o PLD. A CCEE devolve HTTP 403 por
três vias distintas, medido em 06/09/2026. Nada aqui é afirmação sobre preço de
liquidação.

\campo{Grão de subsistema} O ONS não publica CMO nodal. O comparador direto da
literatura, nó a nó, não é reproduzível com este dado: objeto de comparação
inexistente, não validação pendente.

\campo{Uma configuração e uma semente} O intervalo cobre variação entre dias do
teste. Não cobre variância de ajuste, de semente, nem da data de corte entre
treino e teste. A ablação dá uma medida parcial dessa variância.

\campo{Custo} \num{864}~s de ajuste no resultado principal, 15 recalibrações.
Recalibração diária, que a literatura de preço pede, exigiria 444 e não foi
testada.

\conclusao{A um dia, nenhum modelo se distingue da ingênua sazonal em MAE; a
dois dias o modelo ganha 16,31\% com intervalo inteiro acima de zero. O trunfo
da ingênua é acertar o piso do CMO exatamente, e ele é perecível.}

\section*{Artefatos}

% Nome de arquivo em monoespacada nao quebra sozinho, e justificar a coluna
% estreita da ficha deixa hbox subcheia em cada linha. Ragged aqui, e so aqui.
\begin{ficha}
  \item[Código] \raggedright \texttt{experiments/E-preco-cmo/}, cinco passos
                numerados mais \texttt{\_comum.py}, \texttt{\_painel.py} e
                \texttt{\_modelos\_comum.py}.
  \item[Saídas] \raggedright \texttt{data/interim/E-preco-cmo/}:
                \texttt{painel.parquet}, \texttt{janelas.json},
                \texttt{controles.json}, \texttt{previsoes.parquet},
                \texttt{resultado.json}, \texttt{decomposicao.json},
                \texttt{ablacao.json}.
  \item[Registro] \raggedright \texttt{TAREFAS.txt}, entrada F11, e
                  \texttt{report/tarefas-dados.tex}, tarefa~11.
\end{ficha}

\end{document}
